\begin{problem}[Term independent of $x$]
Find the term independent of $x$ in the expansion of
\[
\left(11x-\frac{4}{x^2}\right)^{15}.
\]
\end{problem}

\begin{hint}
Use the general term in the binomial expansion. Then make the power of $x$ equal to $0$.
\end{hint}

\begin{solution}
The general term is
\[
T_{r+1}=\binom{15}{r}(11x)^{15-r}\left(-\frac{4}{x^2}\right)^r.
\]
So the power of $x$ is
\[
x^{15-r}x^{-2r}=x^{15-3r}.
\]
For the term independent of $x$, we need
\[
15-3r=0 \quad \Rightarrow \quad r=5.
\]
Hence the required term is
\begin{align*}
\binom{15}{5}(11x)^{10}\left(-\frac{4}{x^2}\right)^5
&= \binom{15}{5}\cdot 11^{10}\cdot (-4)^5 \\
&= 3003 \cdot 11^{10}\cdot (-1024).
\end{align*}
Therefore, the term independent of $x$ is
\[
\boxed{-3003\cdot 1024\cdot 11^{10}}.
\]
\end{solution}

\begin{takeaways}
\item For ``term independent of $x$'' questions, write the general term first.
\item The key step is solving the exponent equation that makes the power of $x$ equal to $0$.
\end{takeaways}

\begin{problem}[Five-letter words from \texttt{Mamungkukumpurangkuntjunya}]
How many $5$-letter ``words'' can be made from the word \texttt{Mamungkukumpurangkuntjunya}?

\smallskip
\noindent\textit{Helpful note:} To save time, you may use the repeated-letter frequencies
\[
u^6,\ n^4,\ m^3,\ a^3,\ k^3,\ g^2,
\]
and note that the remaining letters each occur once.
\end{problem}

\begin{hint}
Treat this as an arrangement problem with repeated letters. Split into cases according to the repetition pattern of the $5$ letters.
\end{hint}

\begin{solution}
The word \texttt{Mamungkukumpurangkuntjunya} contains the letter types
\[
a,g,j,k,m,n,p,r,t,u,y
\]
with frequencies
\[
u^6,\ n^4,\ m^3,\ a^3,\ k^3,\ g^2,
\]
and $p,r,t,j,y$ each occurring once.

We count by repetition pattern.

\textbf{1. All $5$ letters distinct:}
\[
\binom{11}{5}5! = 55440.
\]

\textbf{2. One pair and $3$ singles:}
\[
6\cdot \binom{10}{3}\cdot \frac{5!}{2!}=43200.
\]

\textbf{3. Two pairs and $1$ single:}
\[
\binom{6}{2}\cdot 9 \cdot \frac{5!}{2!2!}=4050.
\]

\textbf{4. One triple and $2$ singles:}
\[
5\cdot \binom{10}{2}\cdot \frac{5!}{3!}=4500.
\]

\textbf{5. One triple and one pair:}
\[
5\cdot 5\cdot \frac{5!}{3!2!}=250.
\]

\textbf{6. Four of one letter and one single:}
\[
2\cdot 10\cdot \frac{5!}{4!}=100.
\]

\textbf{7. Five of one letter:}
\[
1.
\]

Therefore, the total number of $5$-letter words is
\[
55440+43200+4050+4500+250+100+1=107541.
\]

Hence the answer is $\boxed{107541}$.
\end{solution}

\begin{takeaways}
\item When letters repeat, count by repetition pattern such as $(2,1,1,1)$ or $(3,2)$.
\item After choosing the letters, use multinomial arrangements such as $\frac{5!}{2!}$ or $\frac{5!}{3!2!}$.
\end{takeaways}

\begin{problem}[Constant term from a product of two expansions]
Find the term independent of $x$ in the expansion of
\[
\left(x+\frac{1}{x}\right)^5\left(2x^2-\frac{3}{x}\right)^6.
\]
\end{problem}

\begin{hint}
Write the general term from each bracket, then find which pair of terms gives a total power of $x$ equal to $0$.
\end{hint}

\begin{solution}
From
\[
\left(x+\frac{1}{x}\right)^5,
\]
the general term is
\[
\binom{5}{r}x^{5-r}\left(\frac{1}{x}\right)^r=\binom{5}{r}x^{5-2r}.
\]

From
\[
\left(2x^2-\frac{3}{x}\right)^6,
\]
the general term is
\[
\binom{6}{s}(2x^2)^{6-s}\left(-\frac{3}{x}\right)^s
=\binom{6}{s}2^{6-s}(-3)^s x^{12-3s}.
\]

So in the product, the power of $x$ is
\[
(5-2r)+(12-3s)=17-2r-3s.
\]
For a constant term, we need
\[
17-2r-3s=0.
\]

With $0\le r\le 5$ and $0\le s\le 6$, the only solution is
\[
r=1,\quad s=5.
\]

So the constant term is
\begin{align*}
\binom{5}{1}\binom{6}{5}\cdot 1 \cdot 2^{1}(-3)^5
&= 5\cdot 6\cdot 2\cdot (-243) \\
&= -14580.
\end{align*}

Hence the term independent of $x$ is
\[
\boxed{-14580}.
\]
\end{solution}

\begin{takeaways}
\item In products of expansions, add the exponents from the chosen terms.
\item The constant term occurs when the total exponent of $x$ is $0$.
\item After finding the exponent condition, substitute the matching values back into the coefficients.
\end{takeaways}

\begin{problem}[Adults and children at the cinema]
Three adults and five children go to the cinema and are seated next to each other in a row of eight seats.

How many ways can these eight people sit so that at least two of the adults sit next to each other?
\end{problem}

\begin{hint}
Count all seatings first, then subtract the seatings in which no two adults sit together.
\end{hint}

\begin{solution}
There are
\[
8!
\]
ways to seat the $8$ distinct people in a row.

Now count the arrangements in which no two adults sit together.

First arrange the $5$ children:
\[
5!
\]
ways.

This creates $6$ gaps where adults could be placed:
\[
\_ \ C \ \_ \ C \ \_ \ C \ \_ \ C \ \_ \ C \ \_ 
\]
Choose $3$ of these $6$ gaps for the adults:
\[
\binom{6}{3}
\]
ways.

Then arrange the $3$ adults in those chosen gaps:
\[
3!
\]
ways.

So the number of seatings with no two adults together is
\[
5!\binom{6}{3}3! = 120 \cdot 20 \cdot 6 = 14400.
\]

Therefore, the number of seatings in which at least two adults sit together is
\[
8! - 14400 = 40320 - 14400 = 25920.
\]

Hence the required number of ways is
\[
\boxed{25920}.
\]
\end{solution}

\begin{takeaways}
\item ``At least'' often suggests complementary counting.
\item To keep certain people apart, arrange the others first and then place people into the gaps.
\item Gap methods are especially useful for adjacency restrictions in line arrangements.
\end{takeaways}

\begin{problem}[Lines determined by points on a line and a circle]
There are $11$ points on the plane.

\begin{itemize}[leftmargin=*]
\item $4$ of them lie on the line $y=0$
\item $7$ of them lie on the circle
\[
x^2+(y-2)^2=1.
\]
\end{itemize}

% Assume that, apart from the $4$ points on $y=0$, no other three of the $11$ points are collinear.

How many different lines can be formed that pass through at least two of the $11$ points?

\begin{center}
\begin{tikzpicture}[scale=1.2]
    \draw[->] (-3.2,0) -- (3.2,0) node[right] {$x$};
    \draw[->] (0,-0.6) -- (0,3.6) node[above] {$y$};
    \draw[dashed] (-3,0) -- (3,0) node[right] {$y=0$};
    \draw (0,2) circle (1);
    \node[right] at (1,2.15) {$x^2+(y-2)^2=1$};

    % Four points on the line
    \foreach \x in {-2.4,-1.0,0.9,2.3}
        \filldraw (\x,0) circle (1.5pt);

    % Seven points on the circle
    \foreach \ang in {90,35,-10,-55,-120,165,130}
        \filldraw ({cos(\ang)},{2+sin(\ang)}) circle (1.5pt);
\end{tikzpicture}
\end{center}
\end{problem}

\begin{hint}
Start with all pairs of points, then correct for the overcount caused by the $4$ collinear points on $y=0$.
\end{hint}

\begin{solution}
If no three points were collinear, the number of lines would be
\[
\binom{11}{2}=55.
\]

However, the $4$ points on the line $y=0$ are collinear. The pairs among these $4$ points are
\[
\binom{4}{2}=6,
\]
but they produce only \emph{one} line.

So among these $6$ counted pairs, we have overcounted by
\[
6-1=5.
\]

Therefore, the number of distinct lines is
\[
\binom{11}{2}-5 = 55-5 = 50.
\]

Hence the required number of lines is
\[
\boxed{50}.
\]
\end{solution}

\begin{takeaways}
\item The usual starting point is $\binom{n}{2}$, because each pair of points determines a line.
\item If several points are collinear, those pairs all represent the same line and must be adjusted.
\item A block of $k$ collinear points contributes $\binom{k}{2}$ pairs but only $1$ line.
\end{takeaways}

\begin{problem}[Triangles from points on a circle]
Seven points $P_1,P_2,\dots,P_7$ are arranged in order around a circle, as shown below.

\begin{center}
\begin{tikzpicture}[scale=1.3]
    \draw (0,0) circle (2);
    \foreach \i/\ang in {
        1/90,
        2/38.57,
        3/-12.86,
        4/-64.29,
        5/-115.71,
        6/-167.14,
        7/141.43
    }{
        \filldraw ({2*cos(\ang)},{2*sin(\ang)}) circle (1.6pt);
        \node at ({2.35*cos(\ang)},{2.35*sin(\ang)}) {$P_{\i}$};
    }
\end{tikzpicture}
\end{center}

\begin{enumerate}[label=\roman*)]
\item How many triangles can be drawn using these points as vertices?
\item How many pairs of triangles can be drawn, where the vertices of each triangle are distinct points?
\end{enumerate}
\end{problem}

\begin{hint}
For part (i), choose any $3$ of the $7$ points. For part (ii), choose $6$ points, then split them into two groups of $3$.
\end{hint}

\begin{solution}
\textbf{i)} Any $3$ of the $7$ points form a triangle, so the number of triangles is
\[
\binom{7}{3}=35.
\]

\textbf{ii)} First choose the $6$ points that will be used:
\[
\binom{7}{6}=7.
\]
Now split those $6$ points into two groups of $3$ to make two triangles:
\[
\binom{6}{3}=20.
\]
This counts each pair of triangles twice, because choosing Triangle A first or Triangle B first gives the same pair. So divide by $2$:
\[
\frac{\binom{7}{6}\binom{6}{3}}{2}
=\frac{7\cdot 20}{2}
=70.
\]

Hence the answers are
\[
\boxed{35 \text{ and } 70}.
\]
\end{solution}

\begin{takeaways}
\item On a circle, any choice of $3$ distinct points gives a triangle.
\item When counting pairs of groups, check whether the order of the groups matters.
\item If the pair is unordered, divide by $2$ after choosing the first and second groups.
\end{takeaways}
